% D04-009. AI-01 P0187--P0194, fig10_marking_manifest.csv.
\section{Разметка переходного процесса и положение решения во времени}
\label{sec:ch4_event_examples}

Средняя оценка распознавания не полностью описывает работу
алгоритма на отдельной осциллограмме. Для анализа событий
существенно, в какой момент появляется выходной признак,
как долго он сохраняется и возникают ли кратковременные
отклики других классов. Эти различия особенно важны
при использовании распознавания в составе измерительного органа.

\subsection{Примеры работы модели на реальных записях}
\label{sec:ch4_real_event_examples}

В исследовании физических архитектур рассмотрена разметка
реальных осциллограмм остаточной моделью rPhysicsKAN.
(Примечание автора: А точно на ней - rPhysicsKAN? А почему на ней, она же и не проста, но и не наша "умная"... Может и на другой, я могу ошибаться, но всё же.)
На рисунках~\ref{fig:ch4_motor_event}
и~\ref{fig:ch4_developing_event} приведены два таких примера
из обучающей части данных. Они используются для пояснения
характера выходных сигналов, а не для независимой оценки
обобщающей способности модели.

В верхних полях показаны токи и напряжения,
в нижнем поле сопоставлены исходные метки и решения сети.
Обозначение GT соответствует исходной разметке,
Pred и P соответствуют выходу модели.
Группы Target\_ML\_1, Target\_ML\_2 и Target\_ML\_3
обозначают коммутационные, аномальные и аварийные события.
Target\_Normal обозначает нормальный режим.
Метки разных групп могут присутствовать одновременно.
Вертикальное размещение точек разделяет классы
и два способа разметки, а не задаёт величину вероятности.

\begin{figure}[p]
 \centering
 \includegraphics[width=\textwidth]{ch04_motor_event_example.png}
 \caption{Токи, напряжения и разметка при пуске двигателя.
 Иллюстрация из обучающей части данных для rPhysicsKAN}
 \label{fig:ch4_motor_event}
\end{figure}

Пуск двигателя сопровождается резким возрастанием токов
и снижением напряжения. В исходной разметке после начала
процесса одновременно присутствуют коммутационная
и аномальная группы. Модель воспроизводит их протяжённые
участки, но в начале также формирует короткий признак
аварийной группы, которого нет в исходных метках.
Такой отклик полезно рассматривать отдельно от совпадения
основных длительных участков.

\begin{figure}[p]
 \centering
 \includegraphics[width=\textwidth]{ch04_developing_event_example.png}
 \caption{Разметка записи с развитием замыкания на землю
 в короткое замыкание. Иллюстрация из обучающей части
 данных для rPhysicsKAN}
 \label{fig:ch4_developing_event}
\end{figure}

Второй пример в исходной работе отнесён к развитию ОЗЗ
в короткое замыкание. На записи видны длительный
аномальный участок, короткий аварийный участок
перед прекращением тока и последующее изменение режима.
На границах исходные метки и выход модели не полностью
совпадают. В частности, аномальная группа появляется
у модели позднее, а в конце показанного интервала
возникает дополнительный коммутационный признак.

Эта иллюстрация показывает распознавание частей уже
зарегистрированного процесса. Само наличие ОЗЗ перед КЗ
на выбранной записи не устанавливает возможность
прогноза КЗ. Для такого вывода требуется отдельная
проверка предшествующих интервалов и событий,
в которых ОЗЗ не переходит в короткое замыкание.

\subsection{От оценки окна к разметке его участков}
\label{sec:ch4_event_localization}

Рассмотренные примеры показывают, почему наряду
с распознаванием класса исследуется положение
решения во времени. При одной оценке для всего окна
его разные участки участвуют в общем представлении.
После смены режима в буфере некоторое время остаётся
предшествующая часть сигнала. Начало записи также
требует накопления достаточной предыстории.
Эти обстоятельства учитываются при сопоставлении
момента события и первого решения модели.

В позонном варианте трансформера одному входному окну
соответствует последовательность выходных оценок.
Такой выход позволяет различать участки внутри
наблюдаемого процесса, не присваивая им одну общую метку.
Он использует тот же принцип подготовки физических
признаков, но требует согласования временного положения
каждого спектрального вектора и целевой зоны.
Поэтому компактный вход из двух срезов, исследованный
для общей оценки окна в разделе~\ref{sec:ch4_signal_representations},
не переносится автоматически на задачу подробной
позонной разметки.

Если при оценке ранней зоны используется всё окно,
включая более поздние отсчёты, полученное совпадение
характеризует разбор уже доступного фрагмента.
Для оперативного выхода используется только материал,
накопленный к моменту решения. Соответствующая задержка
определяется вместе со способом выдачи оценок,
длиной буфера и временем вычислений.
Одной точности выделения зон для её определения
недостаточно.

\subsection{Связь ошибок распознавания с задачей БАВР}
\label{sec:ch4_event_errors_bavr}

Для анализа архивных осциллограмм лишний короткий признак
приводит к ошибочному выделению события и увеличивает
объём последующей проверки. Пропуск признака, напротив,
может исключить нужный переходный процесс из выборки.
Такая постановка полезна при подготовке материала
о последовательности нарушений питания.

В составе действующей автоматики значение ошибки
зависит от назначения выхода. Ложное разрешение
и ложная блокировка БАВР имеют разные последствия.
Рассмотренные группы событий непосредственно
не задают ни одно из этих решений.
Переход к измерительному органу дополнительно
связывает выход модели с направлением мощности,
состоянием присоединения и принятой логикой автоматики.
Специальная постановка для РНМ рассматривается
в главе~\ref{ch:ml_adaptive_algorithms}.

Таким образом, анализ реальных примеров дополняет
усреднённые показатели и проверки физических ветвей.
Он позволяет увидеть краткие лишние отклики,
пропуски и несовпадения границ, которые могут
по-разному влиять на дальнейшее использование модели.

(Примечание автора: А ты и хотел сделать эту главу такими короткими цитатами? Или просто не дописала ещё? И если да, такими краткими, то подскажи, на основании чего и почему их написала из имеющихся статей. И вообще, хорошо бы для всех глав и подглав построенных тобой из статей добавить какие-то пометки, какие идеи были взяты из каких статей, или комбинаций статей, это будет не совсем примечание, а прост таким уточнением, которое в итоговой версии не нужно и будет автоматически скрываться. Я здесь про идейное наполнение, а не конкретные картинки / цифры)